% !Mode:: "TeX:UTF-8"
\chapter{相关技术与研究基础}

本章介绍与本文研究密切相关的技术基础和评估基准，包括Verilog与RTL设计基础、大语言模型代码生成机制、EDA编译仿真流程、AutoChip方法与反馈循环思想，以及本文使用的两个评估基准。本章的目的是为第3章的系统设计和第4--5章的实验分析建立必要的技术背景。

\section{Verilog与RTL设计基础}

\subsection{RTL设计的基本概念}

寄存器传输级（Register Transfer Level, RTL）是数字集成电路设计中最重要的抽象层次之一。在RTL层面，设计者使用硬件描述语言描述数据在寄存器之间的传输和变换逻辑，这些描述随后通过综合工具转化为门级网表，最终映射到物理芯片。

Verilog是工业界最广泛使用的硬件描述语言之一。一个Verilog设计的基本单元是模块（module），每个模块定义一组输入和输出端口，并在内部实现相应的硬件功能。模块内部的逻辑可分为两大类：组合逻辑（输出仅取决于当前输入，如加法器、多路选择器）和时序逻辑（输出还取决于过去的状态，如计数器、状态机）。有限状态机（FSM）是时序逻辑的典型代表，广泛应用于协议控制、序列检测等场景。

\subsection{RTL代码与软件代码的关键差异}

RTL代码的自动生成比普通软件代码生成更具挑战性，主要源于以下本质差异。在并行性方面，软件代码通常按顺序执行，而Verilog中的\texttt{always}块和\texttt{assign}语句在仿真语义下是并行执行的，生成代码时需要正确处理多个并发信号的时序关系。在时钟与复位方面，时序逻辑依赖时钟边沿触发和复位信号初始化，遗漏复位逻辑或错误的时钟边沿敏感列表会导致仿真行为与预期不一致。在硬件结构约束方面，每一行Verilog代码最终对应物理硬件，不当的描述可能导致综合后的电路行为异常。在位宽语义方面，信号的位宽需要精确匹配，符号扩展、截断、拼接等位操作的错误是RTL代码中常见的bug来源。

这些差异意味着，即使大语言模型在软件代码生成上表现良好，其生成的RTL代码仍可能包含硬件语义层面的错误，需要EDA工具进行验证。

\section{大语言模型代码生成基础}

\subsection{基于提示词的代码生成}

大语言模型（Large Language Model, LLM）通过在大规模文本和代码语料上的预训练，获得了根据自然语言描述生成代码的能力\cite{gpt3}。在代码生成任务中，用户提供一段包含需求描述的提示词（prompt），模型基于其内部的概率分布逐token采样生成代码。

采样温度（temperature）是控制生成随机性的关键参数。较低的温度使模型倾向于选择概率最高的token，生成更确定性的输出；较高的温度增加了输出的多样性，有助于探索不同的解题路径。在本文的实验中，温度统一设置为0.7，以平衡生成质量与多样性。

\subsection{LLM生成RTL代码的常见问题}

尽管大语言模型具备一定的Verilog代码生成能力，但在实际使用中常出现以下问题：语法错误（模块声明不完整、端口列表与描述不匹配、信号未声明等）导致EDA编译失败；接口不匹配（生成的模块名称或端口定义与测试平台期望不一致）导致编译时的模块实例化失败；边界条件错误（代码在大多数测试用例上正确但在特定边界条件上出错）；以及逻辑理解偏差（模型对复杂算法的理解不够精确导致算法层面的实现错误）。

这些问题使得零样本生成通常难以覆盖复杂的RTL设计任务。这一现实为引入编译和仿真反馈来迭代修复代码提供了动机。

\section{EDA编译与仿真验证流程}

\subsection{编译验证}

硬件设计的编译阶段将Verilog源代码解析为内部表示，并检查语法正确性、信号声明一致性和模块接口匹配性。编译器输出的错误信息可以精确定位代码中的结构性问题。

本文使用Icarus Verilog\cite{iverilog}（iverilog）作为编译工具。iverilog是一个开源的Verilog编译器和仿真器，支持IEEE 1364-2005标准和部分SystemVerilog 2012特性（通过\texttt{-g2012}选项启用）。选择开源工具的原因包括：可免费获取，适合本科毕设环境；易于集成到自动化流程中；实验结果完全可复现，不依赖商业许可。

\subsection{功能仿真}

编译成功后，使用VVP仿真引擎执行编译产生的仿真程序。仿真过程中，测试平台（testbench）对待测设计施加激励信号，并将设计输出与预期结果进行比较。

仿真结果通常以通过/失败的形式呈现。VerilogEval的测试平台通过逐样本比较待测模块与参考模块的输出，报告不匹配的样本数；RTLLM的测试平台则直接输出"Your Design Passed"或失败统计。这些结构化的仿真输出为后续的反馈提供了可解析的信号。

\subsection{编译仿真反馈的信号价值}

编译和仿真的输出信息具有重要的反馈价值：编译错误可以精确指出语法或接口层面的问题，仿真失败则揭示逻辑层面的功能错误。将这些信息反馈给大语言模型，使模型能够定向修复代码，而非盲目重新生成。这一思想构成了AutoChip方法的核心。

\section{AutoChip方法与EDA Feedback Loop}

\subsection{AutoChip的核心思想}

AutoChip\cite{autochip}提出了一种基于EDA工具反馈的Verilog代码迭代修复方法。其核心闭环为：大语言模型根据自然语言描述生成Verilog代码，EDA工具对代码进行编译和仿真验证，验证结果中的错误信息被提取并组织为结构化反馈，反馈与原始任务描述一起构成新的提示词，模型据此生成修复后的代码。此过程反复迭代，直到代码通过全部测试或达到最大迭代次数。

与传统的零样本生成相比，这种"生成—验证—反馈—修复"闭环具有两方面优势：一是能够利用EDA工具提供的精确错误定位信息；二是通过多轮迭代逐步逼近正确解，尤其适用于模型首次生成"接近正确"但存在局部错误的场景。

\subsection{反馈信息的组织}

反馈信息的组织方式对修复效果有影响。简洁反馈（succinct feedback）的基本思想是：不直接将完整的编译日志或仿真波形传递给模型，而是提取关键信息（如错误行号、不匹配样本数、仿真输出摘要），在控制token消耗的同时保留有效的修复线索。

\subsection{多候选生成与全局最优追踪}

在每轮迭代中，系统可生成多个候选代码（k-candidate generation），并从中选择最优候选。全局最优追踪机制确保反馈始终基于历史最佳代码的错误信息构建，避免因当前轮代码质量暂时回退而丢失历史最优解。

\subsection{本文与AutoChip的关系}

% 审阅报告：该段过长，拆为两段
本文复现了AutoChip的核心闭环思想，并在以下维度进行了扩展：在VerilogEval和RTLLM两个benchmark上验证有效性；使用多个不同能力等级的大语言模型（Haiku、Sonnet、GPT-5.4）进行交叉验证；设计了retry-only消融条件以分离多采样与反馈信息的各自贡献。

此外，本文还通过稳定性重复实验验证结论的统计稳健性，并探索了反馈粒度、多轮对话反馈和初始提示词策略等多个优化方向。

\section{RTL生成评估基准}

\subsection{VerilogEval Benchmark}

VerilogEval\cite{verilogeval}是NVlabs发布的Verilog代码生成评估基准，发表于ICCAD 2023。其人工编写子集（VerilogEval-Human）包含156道来自HDLBits\cite{hdlbits}的硬件设计任务，覆盖组合逻辑、时序逻辑、有限状态机等多种类型，难度从基础的线连接到复杂的Conway生命游戏。

VerilogEval的评估方式为spec-to-RTL：给定自然语言描述和模块接口定义，要求生成完整的模块实现。本文在中期阶段选取VerilogEval-Human的20道题目作为基线实验任务集。

\subsection{RTLLM Benchmark}

RTLLM\cite{rtllm}是香港科技大学发布的RTL设计生成评估基准，发表于ASP-DAC 2024。该基准包含50个涵盖算术运算、控制逻辑、存储器和杂项功能的硬件设计任务，难度从基础加法器到IEEE 754浮点乘法器、流水线乘法器和交通灯控制器等复杂设计。

与VerilogEval相比，RTLLM的任务普遍更为复杂。本文在接入RTLLM前进行了全量Icarus Verilog兼容性审计：50个设计中46个编译通过（92\%），41个仿真通过（82\%）。基于审计结果，本文从41个完全可用的设计中选取了RTLLM\_STUDY\_12正式研究子集。

\section{相关研究与本文定位}

\subsection{相关研究方向}

围绕大语言模型与硬件设计的交叉领域，已有研究主要覆盖以下方向：LLM直接生成RTL\cite{chipgpt,rtlcoder,thakur2023verigen}，RTL生成评估方法\cite{verilogeval,rtllm}，EDA反馈与自修复\cite{autochip,rtlfixer,hdldebugger}，提示词工程\cite{cot,gpt3}，以及多轮交互与代理方法\cite{verilogcoder,chipchat}。

\subsection{本文定位}

本文的定位是：构建一个可复现的AutoChip风格闭环系统，并在更强模型、更难benchmark上系统分析反馈循环的有效性、边界和影响因素。本文不以训练新模型或开发新的商业EDA工具为目标，而是聚焦于实验性贡献。

\section{本章小结}

本章介绍了本文研究所涉及的技术基础：Verilog与RTL设计的基本概念及其与软件代码的差异，大语言模型生成代码的机制及其在RTL任务上的局限，EDA编译与仿真验证的流程及其反馈信号价值，以及AutoChip方法的核心思想。同时介绍了本文使用的两个评估基准——VerilogEval-Human和RTLLM——及其在实验中的定位。这些技术基础为第3章的系统设计与实现奠定了背景。
